\chapter{Introduction}
\label{chap:introduction}

\section{Motivation}

Large language models generate text by repeatedly evaluating a trained
model on a growing sequence of tokens.  Providing useful inference
performance therefore requires both numerical computation and timely
access to substantial amounts of data.  The model's learned parameters
remain available across requests, while attention state grows with the
histories of the sequences being processed.  Temporary results and
runtime workspaces add further demands.  Model size alone consequently
does not describe the memory required to serve a workload.

CPUs provide a relevant setting for studying these demands.  Their
inference performance depends on the numerical representation of the
model, the kernels that execute it, the available parallelism, and the
memory system supplying their operands~\cite{na2024llmcpu}.  A processor
can have considerable arithmetic capability while an application spends
much of its time waiting for data.  Conversely, increasing memory
bandwidth is useful only when the computation can exploit it.

Systems combining high-bandwidth memory (HBM) and conventional DDR
memory make this relationship particularly important.  They provide
memory resources with different bandwidth, capacity and locality
characteristics.  Depending on the hardware configuration, HBM may be
exposed as separately addressable memory or used as a hardware-managed
cache for DDR~\cite{intelmaxoverview}.  Understanding the benefit of such
systems requires considering both the data retained by inference and
the traffic generated while using it.

\section{Objective and scope}

This thesis investigates the memory requirements and performance of
CPU-based LLM inference on systems combining HBM and DDR, considering
the storage and access behavior of model weights and KV state.

The central problem is the relationship between inference execution and
the memory resources available within a compute node.  Weights and
cached keys and values have different origins, lifetimes and access
patterns.  The former encode the trained model; the latter retain
information computed for a sequence's attention history.  Both must be
considered alongside runtime buffers when reasoning about capacity,
data movement and performance.  Prompt length, generation length,
concurrency and execution batching can change this balance even when
the model remains fixed.

The scope is inference with trained transformer language models, with
particular attention to autoregressive generation on CPUs.  Training
and distributed inference across compute nodes are outside the scope
of the memory-system analysis developed here.  The experimental setting
is the CRESCO8 CPU platform described in
\chapref{chap:experimental-platform}.  Architecture and runtime
configuration are relevant because they determine which storage and
execution mechanisms an experiment can actually exercise.

The investigation considers weight storage, KV storage and their
interaction under declared workloads.  Allocation, placement and
scheduling affect the resources available to computation; their effects
must be assessed together with any preparation or data-movement costs.
Comparisons require explicit memory configurations, workload definitions
and timing boundaries, so that observed performance can be related to
the data being stored and accessed.

\section{Thesis organization}

\todo[inline,size=\normalsize]{\textbf{TODO:} Write the thesis organization
once the thesis is complete, describing the final chapters and appendices
in their actual order.}
